% Compile twice: lualatex 30005916.tex
\documentclass[11pt]{article}
\usepackage[margin=25mm]{geometry}
\usepackage{fontspec,unicode-math}
\setmainfont{Latin Modern Roman}
\setmathfont{Latin Modern Math}
\usepackage{amsmath,amssymb,mathtools}
\usepackage{xurl,hyperref}
\hypersetup{colorlinks=true,urlcolor=blue,linkcolor=blue}
\setcounter{secnumdepth}{0}
\setlength{\parindent}{0pt}
\setlength{\parskip}{5pt}
\setlength{\emergencystretch}{3em}
\newcommand{\sref}[2]{\hyperlink{source:#1:#2}{[S#2]}}
\newcommand{\cataloguescope}[1]{\paragraph{Source question.}#1}
\newcommand{\floor}[1]{\left\lfloor#1\right\rfloor}
\begin{document}
% UnsolvedMath ID 30005916; source catalogue code OWR-14298371-007
\setcounter{section}{30005915}
\section{Unconditionally Ruling Out Heuristica}\label{30005916}
{\small\sffamily UnsolvedMath ID 30005916 · Computer Science}
\subsection{Definitions and mathematical statement}

\paragraph{Shared notation.}Standard number systems are denoted by $\mathbb N,\mathbb Z,\mathbb Q,\mathbb R,\mathbb C$; index ranges, computational models and specialized notation follow the cited problem source and the definitions below.


Heuristica is the hypothetical world where worst-case NP hardness exists but NP problems have efficient average-case algorithms under the designated efficiently samplable input distributions. Ruling it out requires the source\textquotesingle s exact notion of average-case easiness.
\paragraph{Statement recorded in the supplied source.}
Can the complexity-theoretic world Heuristica be ruled out unconditionally? \sref{30005916}{2}
\paragraph{Scope qualification.}
The precise distribution family and allowed average-case error are source-dependent; no unconditional separation is claimed.
\paragraph{Source status.}
The supplied catalogue labels this record open. Its source-status metadata was checked against the supplied snapshot on 2026-10-11. This import does not independently establish current open status or mathematical correctness. The catalogue literature assessment is dated 2026-08-22.
\subsection{Short English statement}
Exclude the stated worst-case-hard but average-case-easy complexity world without assuming a hardness hypothesis.
\subsection{Sources}
\begin{itemize}
\item[\hypertarget{source:30005916:1}{S1}] Ulam AI and the UnsolvedMath Contributors, UnsolvedMath: A Curated Collection of Open Mathematics Problems, record 30005916 (OWR-14298371-007). Catalogue curation and metadata: CC BY 4.0. The supplied source question is retained; definitions and the English summary are editorial additions.
\url{https://huggingface.co/datasets/ulamai/UnsolvedMath}.
\item[\hypertarget{source:30005916:2}{S2}] Complexity Theory (2024), Oberwolfach Reports, DOI 10.4171/owr/2024/27, pp. 1517--1562
\url{https://doi.org/10.4171/owr/2024/27}.
\item[\hypertarget{source:30005916:3}{S3}] Hirahara--Nanashima
\url{https://eccc.weizmann.ac.il/report/2024/063/download/}.
\item[\hypertarget{source:30005916:4}{S4}] ECCC 2026
\url{https://eccc.weizmann.ac.il/report/2026/052/}.
\textit{Underlying Oberwolfach Reports material retains its source-specific licence, including CC BY-SA 4.0 where applicable; see THIRD\_PARTY\_NOTICES.md.}
\end{itemize}
\paragraph{Preparation and review.}
Original question attributed to the cited source and UnsolvedMath contributors. Codex (GPT-6; exact serving revision not exposed) assisted with definitions, English summaries, source inspection and formatting. No new human mathematical authorship or independent proof verification is asserted. Canonical ID 30005916 and source code OWR-14298371-007 checked against the complete supplied registry and existing repository statements; known flagged repetitions and obvious identical conjectures excluded. Full mathematical-equivalence/literature review has not been performed. Statement preparation and final review followed contribution prompts C and D; automated checks verify structure and identity, not open status.
\label{end:30005916}
\end{document}
